\section{¿Es el modelo fundacional para robots una disciplina independiente? So far, not yet}

La sección anterior trató los dos cambios de paradigma; esta sección analiza una controversia: ¿debe considerarse el gran modelo fundacional para robots una disciplina independiente de los grandes modelos? La respuesta de Tan Jie (谭杰) es mesurada: so far, not yet. Hoy, la mayor parte de la inteligencia robótica sigue dependiendo de grandes modelos multimodales; lo único que se hace es complementarlos con salida de robot action, datos de acciones y fine-tuning por etapas. Si en el futuro aparecen un nuevo data format, modelos del mundo o cuellos de botella de control, podría convertirse en una disciplina más independiente, pero por ahora no ha habido un cambio cualitativo.

\begin{figure}[H]
\centering
\includegraphics[width=0.96\textwidth]{vla-action-gap.png}
\caption{VLA completa la salida de Action: para controlar un robot, un modelo multimodal tiene que aprender a producir acciones. Diagrama conceptual propio, elaborado a partir de la conversación de 00:13:06--00:23:44.}
\end{figure}

\begin{knowledgebox}{Lectura de la figura: de VLM a VLA, lo decisivo son los datos de acciones}
Vision y Language entran al modelo multimodal, que originalmente produce texto y planes; un robot necesita acciones, por lo que se incorpora Action Data para que el modelo pase de VLM a VLA. Aquí las acciones no son texto abstracto, sino articulaciones, trayectorias, efectores finales y señales de control del robot.
\end{knowledgebox}

\subsection{VLA, generalización y tasa de éxito}

Esta subsección lleva la «salida de acciones» al terreno de la utilidad práctica. La importancia de VLA no se reduce a añadir action al nombre: consiste en que el robot pase efectivamente del lenguaje y la visión a acciones ejecutables.

VLA significa Vision-Language-Action, es decir, un modelo de visión, lenguaje y acción. Busca reunir «ver el entorno», «entender la instrucción» y «producir la acción» en un mismo modelo o sistema. Los modelos actuales alcanzan tasas de éxito muy altas en tareas sencillas como pick-and-place, pero en la manipulación fina, por ejemplo subir un cierre, sujetar con precisión o controlar la orientación, la tasa de éxito puede ser de apenas 30 o 40\%. Esa cifra representa un avance en un video de investigación, pero no sirve en la vida real.

\begin{warningbox}{La tasa de éxito en investigación no equivale a la utilidad de un producto}
En el mundo real, una manipulación fina con 30\%--40\% de éxito es prácticamente inutilizable. Los usuarios necesitan sistemas estables, recuperables, explicables y que fallen de forma segura, no solo que completen la tarea una vez en un video.
\end{warningbox}

\subsection{De la idea a la demo y a la implantación}

Tan Jie recurre a una analogía con la conducción autónoma: pasar de una idea a un prototype publicado en un artículo puede tomar de seis meses a un año; del artículo a atreverse a hacer una live demo, uno o dos años; y de la live demo a la implantación real, de cinco a diez años. La robótica es más difícil que la conducción autónoma porque el espacio de acciones es mayor, hay más tareas y la interacción física es más compleja. La conducción autónoma es un escenario vertical con una salida de acciones relativamente limitada y, aun así, tardó más de diez años en acercarse a la implantación.

\begin{knowledgebox}{Tabla de etapas de la investigación a la implantación}
\begin{tabularx}{\textwidth}{p{0.22\textwidth}p{0.32\textwidth}X}
\toprule
Etapa & Producto habitual & Riesgo principal \\
\midrule
Idea & Idea de investigación, recipe de entrenamiento, hipótesis preliminar & Puede cumplirse solo en condiciones muy restringidas. \\
Prototype & Artículo, evaluación fuera de línea, video grabado & Muchos supuestos; los casos fallidos se filtran. \\
Live demo & Demostración en vivo, éxito en tareas limitadas & La estabilidad y la recuperación ante anomalías siguen siendo insuficientes. \\
Pilot & Prueba piloto con pocos clientes o escenarios & Surgen presiones de costo, mantenimiento, seguridad y operación. \\
Production & Despliegue comercial replicable & Requiere hardware confiable, un volante de datos y un sistema de servicio. \\
\bottomrule
\end{tabularx}
\end{knowledgebox}

\begin{importantbox}{La escala temporal de la robótica}
Un avance de investigación y la implantación comercial no son lo mismo. Para pasar de la demo a la producción real, un robot tiene que superar el umbral al mismo tiempo en tasa de éxito de las tareas, confiabilidad del hardware, seguridad, costo, mantenimiento, volante de datos y valor para el escenario.
\end{importantbox}

\subsection{Resumen de la sección}

Por ahora, el modelo fundacional para robots se parece más a una extensión de los grandes modelos multimodales que a una disciplina completamente independiente. Lo que de verdad decidirá si se independiza es si en el futuro hacen falta nuevos formatos de datos, modelos del mundo, paradigmas de control y mecanismos de generalización entre distintos cuerpos robóticos.
